\documentclass[runningheads]{llncs}

\usepackage[T1]{fontenc}
\usepackage{amsmath,amssymb}
\usepackage{graphicx}
\usepackage{booktabs}
\usepackage{xcolor}
\usepackage[hidelinks]{hyperref}

\newcommand{\eps}{\varepsilon}
\newcommand{\E}{\mathbb{E}}
\newcommand{\Prob}{\mathbb{P}}
\DeclareMathOperator{\Lap}{Lap}
\DeclareMathOperator{\Var}{Var}

\begin{document}

\title{Forecasting Privacy-Budget Consumption in\\ Repeated DP-SQL Workloads}
\titlerunning{Forecasting Privacy-Budget Consumption in Repeated DP-SQL Workloads}

\author{Refik Can \"Ozta\c{s}}
\authorrunning{R.\,C.\ \"Ozta\c{s}}
\institute{Hacettepe University, Ankara, Turkey\\
\email{canoztas06@gmail.com}}

\maketitle

\begin{abstract}
Differentially private (DP) SQL systems add calibrated noise to each aggregate query and charge it a slice of a fixed privacy budget; once it is spent, no further query can be answered. They account for queries in isolation, yet real analytical workloads \emph{repeat}---dashboards refresh the same tiles, alerting jobs re-poll the same counts---and an exact repeat can reuse an already-noised answer for free (DP \emph{post-processing}). We give a closed-form model that predicts, from a workload's repetition structure alone and \emph{before any query runs}---or, when the repetition distribution is unknown, from an early prefix---how much budget it will spend. The key quantity is the expected number of \emph{distinct} queries, $\E[u_k]=\sum_i[1-(1-p_i)^k]$; from it we obtain the expected spend, a savings ratio, and limit, monotonicity, concentration, and utility consequences. To our knowledge this is the first such forecast computed in \emph{closed form from public structure alone}---$\{p_i\}$, length $k$, budget $B$---with no query executed and the sensitive data never touched. We instantiate the model in a unit-tested middleware (COUNT/SUM/AVG, GROUP BY, top-$k$); a $4{,}155$-trial \emph{self-consistency} campaign matches the formula within $3\%$ in $21/24$ cells, and---as external evidence---on \emph{Redbench}, $30$ Amazon Redshift--derived traces, forecasting from only the first half of each workload cuts the prediction error by $45\%$ over a plug-in estimate.
\keywords{Differential privacy \and Privacy budget \and SQL workloads \and Unseen-species estimation \and Caching.}
\end{abstract}

\section{Introduction}
\label{sec:intro}

Aggregate statistics leak. Two COUNT queries that differ by one predicate reveal an individual's attribute by differencing; repeated, such queries reconstruct it~\cite{dinur2003revealing}. Differential privacy (DP)~\cite{dwork2006calibrating} is the standard defense: perturb each answer with noise calibrated to its sensitivity and a per-query budget $\eps_q$, and bound the cumulative loss by composition. The budget is a hard, non-renewable resource: once the total reaches a cap $B$, no further query can be answered. A custodian therefore needs to predict \emph{how fast a workload spends its budget}.

Existing DP-SQL systems charge every query a fresh $\eps_q$ and account for them in isolation. Yet analytical workloads are highly repetitive: a monitoring dashboard re-issues the \emph{exact same} tile on a schedule (a canonicalized exact-query repeat, served free), while a drill-down reuses a template with \emph{different} literals (a distinct query, charged afresh). Under exact-repeat caching a repeated answer is free---reusing an already-released noisy value adds no privacy cost, since any function of a DP output is itself DP (\emph{post-processing})---so the budget a length-$k$ workload consumes is governed not by $k$ but by its number of \emph{distinct} queries $u_k$ (Fig.~\ref{fig:concept}). We ask whether $u_k$, and hence the budget, can be \emph{predicted from the workload's repetition structure, in closed form, before the workload runs}. We claim no novelty for exact-repeat caching---it is a known mechanism; the contribution is \emph{forecasting its budget impact before execution}.

\begin{figure}[t]
\centering
\includegraphics[width=\textwidth]{fig_concept}
\caption{Naive composition charges every one of $k$ queries; workload-aware accounting charges only the $u_k$ distinct query species, since exact repeats are served from cache by post-processing. The model predicts $\E[u_k]$, and hence the budget, before execution.}
\label{fig:concept}
\end{figure}

\paragraph{Contributions.}
\begin{enumerate}
\item \textbf{A before-execution budget forecast} (Sec.~\ref{sec:model})---to our knowledge the first that is \emph{closed-form} and reads only public structure $(\{p_i\},k,B)$, never touching the data. It delivers three outputs of different strengths: the \emph{expected} spend $\eps_q\E[u_k]$; with the safe sizing $\eps_q{=}B/m$, a \emph{hard} completion guarantee ($u_k\le m$ in the static regime, i.e.\ no freshness-driven re-noising); and a \emph{per-release} $(\alpha,\beta)$-accuracy verdict (each release individually, not jointly). Five propositions establish it, with two limits recovering the folklore $1-1/k$ caching rule and standard per-query accounting.
\item \textbf{A measured improvement on real traces.} On $30$ Redbench workloads (from Amazon Redshift logs), forecasting the full distinct-query count from only the \emph{first half} via an unseen-species correction cuts the normalized error by $45\%$ ($0.22\!\to\!0.12$) over the plug-in (Sec.~\ref{sec:eval}).
\item \textbf{A middleware and validation} (Secs.~\ref{sec:system}--\ref{sec:eval})---DuckDB/PostgreSQL, $134$ unit tests; the forecast matches within $3\%$ in $21/24$ cells, invariant to data scale and $\eps_q$.
\end{enumerate}
We do \emph{not} claim to outperform deployed DP caches: against them we are \emph{on par} on the savings number---the contribution is the forecast they do not expose.

\section{Related Work}
\label{sec:related}

\paragraph{DP-SQL systems.} PINQ~\cite{mcsherry2009privacy} and Chorus~\cite{johnson2020chorus} expose a fixed per-query $\eps_q$ and account in isolation. PrivateSQL~\cite{kotsogiannis2019privatesql} spends a one-time budget building DP synopses over a representative workload; it fixes that budget up front but emits no closed-form, before-execution prediction for an arbitrary repeated stream.

\paragraph{DP caching and proactive answering.} DP caches make repeated work cheaper at runtime: Turbo~\cite{kostopoulou2023turbo} conserves $1.7$--$15.9\times$ budget by generalizing across overlapping linear queries, and CacheDP~\cite{mazmudar2023cachedp} answers a workload at lower budget under an accuracy contract. Pioneer~\cite{peng2013pioneer} reuses past noisy results but, by its authors' statement, optimizes only the \emph{current} query and leaves future-workload effects to future work. Closest to us is the concurrent, unrefereed LAPRAS~\cite{mundra2026lapras}, also \emph{proactive}: it precomputes an offline batch mechanism on a \emph{predicted} query set before the stream and serves it for free---but it needs a concrete predicted set and \emph{touches the sensitive data} ahead of time, whereas we forecast in closed form from \emph{public} structure $(\{p_i\},k,B)$ alone, the data never touched. These deliver the \emph{savings} our model prices, but reactively or by precomputing on data, and without a structural forecast; our exact-repeat caching is a special case for which we claim no novelty.

\paragraph{Accounting, allocation, unseen species.} Tight composition (RDP~\cite{mironov2017renyi}, zCDP~\cite{bun2016concentrated}) reduces the \emph{rate} of spend, orthogonal to our \emph{count} of releases. \emph{Workload-aware} mechanisms in the established sense (the matrix mechanism~\cite{li2015matrix}, HDMM~\cite{mckenna2018hdmm}) optimize noise for a \emph{known} batch of linear queries by running on the data, and budget \emph{schedulers} (DPF/PrivateKube~\cite{luo2021dpf}, Cohere~\cite{kuchler2024cohere}) treat the budget as a schedulable resource across competing consumers at runtime; neither forecasts a repeated stream's spend in closed form before execution. APEx~\cite{ge2019apex} maps a per-query accuracy target to the minimum $\eps$ at runtime; multi-analyst budget sharing~\cite{pujol2023multianalyst} invokes the coupon-collector problem to bound queries-per-\emph{analyst}, a throughput limit, not a distinct-query forecast. The distinct-count side is classical occupancy/coupon-collecting~\cite{flajolet1992}; predicting the not-yet-seen remainder is the unseen-species problem: Good--Toulmin~\cite{goodtoulmin1956} and its smoothed variant~\cite{orlitsky2016unseen} predict new species from a prefix; this has even been brought under DP for \emph{estimation} (INSPECTRE~\cite{acharya2018inspectre}), but for support coverage, not budget. We connect these tools to DP-\emph{budget} forecasting.

\section{The Model}
\label{sec:model}

The budget-spending unit (``species'') is the \emph{exact} query---a template together with its WHERE literals, i.e.\ the cache key. Let $m$ be the number of distinct exact queries an analyst may issue and $\{p_i\}_{i=1}^m$ a distribution over them; the analyst issues $k$ queries i.i.d., and $u_k$ counts the distinct queries that appear at least once.

\begin{proposition}[Expected distinct queries]\label{prop:eu}
$\E[u_k]=\sum_{i=1}^m\bigl[1-(1-p_i)^k\bigr]$.
\end{proposition}
\noindent\emph{Proof.} With $X_i=\mathbf 1[\text{query }i\text{ appears}]$, $\Prob[X_i{=}0]=(1-p_i)^k$; linearity gives the sum. \hfill$\square$

\smallskip\noindent\emph{In plain terms,} $(1-p_i)^k$ is the chance query $i$ never appears; one minus that is the chance it appears at least once; summing counts the distinct queries---the only ones that cost budget. \emph{Example:} three tiles asked $70/20/10\%$ over $k{=}100$ refreshes give $\E[u_k]\approx3$: you pay for $\approx3$ queries, not $100$.

\begin{proposition}[Budget and savings]\label{prop:sav}
Naive composition spends $k\eps_q$, while workload-aware accounting spends $\E[\eps_{\mathrm{wa}}]=\eps_q\E[u_k]$ in expectation---a savings ratio $S(k)=1-\E[u_k]/k$.
\end{proposition}
\noindent Three limiting cases check the corners. \emph{(A)} deterministic repetition ($p_1{=}1$): $S=1-1/k$, the folklore caching rule. \emph{(B)} uniform with $m\to\infty$: $\E[u_k]\to k$, $S\to0$, recovering naive composition. \emph{(C)} uniform with $k\to\infty$, $m$ fixed: the budget saturates at $m\eps_q$, so a finite-template workload cannot exhaust a budget large enough to cover all $m$ queries.

\begin{proposition}[Skew monotonicity]\label{prop:zipf}
For $p_i\propto i^{-\alpha}$, $S(k)$ is non-decreasing in $\alpha$ for fixed $(k,m)$.
\end{proposition}
\noindent\emph{Proof sketch.} $\E[u_k]=\sum_i\phi(p_i)$ with $\phi(x)=1-(1-x)^k$ concave, hence $\E[u_k]$ is Schur-concave. For $\alpha_2>\alpha_1$ the ratio $p_i(\alpha_2)/p_i(\alpha_1)\propto i^{-(\alpha_2-\alpha_1)}$ is strictly decreasing in $i$, so the two equal-sum, sorted vectors cross exactly once, from above to below ($p_1(\alpha_2)>p_1(\alpha_1)$ is forced by the equal sums and the decreasing ratio); a single $+$-to-$-$ crossing gives partial-sum dominance, i.e.\ $p(\alpha_2)$ majorizes $p(\alpha_1)$, weakly lowering $\E[u_k]$ (strictly for $k\ge2$) and raising $S$ (a numerical scan over $3\le m\le100$, $2\le k\le200$ corroborates). So the more a workload concentrates on a few popular queries, the more it saves.

\begin{proposition}[Concentration]\label{prop:conc}
With $q_i=1-(1-p_i)^k$, the occupancy indicators are negatively associated~\cite{joagdev1983} (so their covariances are $\le0$), giving $\Var[u_k]\le\sum_i q_i(1-q_i)\le m/4=:V$ and $\Prob[|u_k-\E[u_k]|\ge t]\le V/t^2$.
\end{proposition}
\noindent Because $V=O(m)$ not $O(k)$, a single run's realized spend deviates from the forecast only on the $O(\sqrt m)$ scale, independent of $k$---useful for sizing, though a variance/Chebyshev-scale statement, not a high-probability envelope.

\begin{proposition}[Utility at fixed budget]\label{prop:util}
At a fixed total $\eps_{\mathrm{tot}}$, naive affords $\eps_{\mathrm{tot}}/k$ per query and workload-aware $\eps_{\mathrm{tot}}/\E[u_k]$ per distinct release; since $\E[|\eta|]=\Delta f/\eps_q$, the per-query error ratio is $k/\E[u_k]$.
\end{proposition}
\noindent So spreading the budget over the $\E[u_k]$ distinct queries (not all $k$) buys proportionally lower error---in the example, $\approx33\times$ at equal privacy (in expectation; a hard cap is needed when realized $u_k$ exceeds $\E[u_k]$, which is why the static-regime safe sizing is $\eps_q=B/m$). Beyond the mean, the per-release accuracy has a closed tail: for $\eta\sim\Lap(\Delta f/\eps_q)$---where $\Delta f$ is the sensitivity of \emph{that release's} query ($1$ for COUNT, the column bound for SUM/AVG; Sect.~\ref{sec:system}), so the verdict is query-type-specific---$\Prob[|\eta|>\alpha]=e^{-\alpha\eps_q/\Delta f}$, and each release meets $|\text{error}|\le\alpha$ with probability $\ge1-\beta$ iff $\alpha\ge(\Delta f/\eps_q)\ln(1/\beta)$. Under $\eps_q{=}B/m$ the planner emits this \emph{per-release} $(\alpha,\beta)$ verdict at $t{=}0$, before any query runs---reported per release, not as a joint guarantee.

\section{System}
\label{sec:system}

A thin middleware sits between the analyst and a DuckDB/PostgreSQL backend. Each query is parsed, hashed to its exact cache key, and checked against a cache; a hit is returned free (post-processing), a miss is charged $\eps_q$, noised with $\Lap(\Delta f/\eps_q)$, and cached. Sensitivities are $\Delta f=1$ for COUNT and $\Delta f=B_c$ for SUM/AVG, where the public bound $B_c$ is enforced \emph{by construction}: the executed SQL clamps each summed value to $[0,B_c]$ (NULLs preserved), so out-of-domain data cannot break the guarantee. A GROUP BY costs $\eps_q$ (parallel composition); top-$k$ is post-processing of one noised histogram (order/limit applied on the noised values). The parser enforces a strict positive whitelist---window/arithmetic-wrapped aggregates, self-union CTEs, set operations, derived tables, and OFFSET are rejected before any budget is spent.

The model is also prescriptive: a \emph{predictive allocator} estimates $\hat U\approx\E[u_k]$ online from the empirical query distribution and sets $\eps_q=B/\hat U$, capping spend at $B$. A \emph{temporal} extension prices freshness (staleness $\tau$, update rate $\lambda$) into a first-order planning estimate, not a proven bound; note that re-noising charges fresh budget, so the static $u_k\le m$ cap---and with it the hard completion guarantee---does not extend to the temporal setting.

\paragraph{Scope boundary.} The GROUP BY / top-$k$ guarantees assume a \emph{public, fully enumerated} key domain (absent groups released as noised zeros). Private key \emph{discovery} (a stability threshold) is future work; the prototype returns present keys only, which we state plainly rather than hide.

\section{Evaluation}
\label{sec:eval}

\paragraph{Setup.} A $4{,}155$-trial campaign over Zipf workloads (Adult age-bands and TPC-H templates). The main grid ($m{=}7$, $\eps_q{=}1$) crosses Zipf $\alpha\in\{0,0.5,1,1.5,2,3\}$ with $k\in\{10,25,50,100\}$ at $30$ trials per cell; the savings workloads use $m\in\{1,3,5,7\}$ plus an all-unique drill-down; separate sweeps cover $\eps_q\in\{0.1,0.5,1,2\}$ and $k$ up to $500$, on two data scales (TPC-H SF$=$1/10).

\paragraph{Forecast accuracy.} The closed form matches the empirical mean within $<3\%$ in $21/24$ main-grid cells (the three misses are the smallest-$k$ cells, Zipf $\alpha{\ge}2$ at $k{\le}25$, worst $8.6\%$, where a 30-trial mean is noisiest); it is invariant to data scale (6M vs 60M rows agree within $0.034$ in $\E[u_k]$, ${\approx}0.5\%$ of $m{=}7$) and to $\eps_q$ (it cancels), and is \emph{exact} on a constructed mixed workload. This grid is a \emph{self-consistency} check---trials are drawn from the assumed distribution---confirming the implementation matches the model; the deployment-sizing evidence is the Redbench prefix forecast below.

\paragraph{Savings and the no-win case.} Workload-aware accounting answers the same $100$ queries while spending only $\eps_q u_k$ (Table~\ref{tab:sav}): $100\times$ less budget on a repetitive workload, $14$--$33\times$ on parametric/uniform pools (where $u_k$ saturates at $m$), and \emph{exactly} $1\times$ on a genuine drill-down---which the model correctly predicts as $S(k){=}0$ rather than overclaiming.

\begin{table}[t]
\centering
\caption{Budget saved vs.\ naive composition ($k{=}100$, total $\eps{=}100$, $30$ trials). The rows vary the distinct-query \emph{pool size} $m$; by $k{=}100$ the distinct count $u_k$ saturates near $m$, which the model predicts in closed form. (At fixed $m$, higher skew lowers $u_k$ further; Prop.~\ref{prop:zipf}.)}
\label{tab:sav}
\begin{tabular}{lccc}
\toprule
Workload & naive $\eps$ & workload-aware $\eps$ & saved \\
\midrule
Repetitive (1 query)            & 100 & 1.0  & $100\times$ \\
TPC-H returnflag ($m{=}3$)      & 100 & 3.0  & $33\times$ \\
TPC-H priority ($m{=}5$)        & 100 & 5.0  & $20\times$ \\
Uniform / Zipf pool ($m{=}7$)   & 100 & 7.0  & $14.3\times$ \\
Drill-down (all unique)         & 100 & 100  & $1\times$ (no win) \\
\bottomrule
\end{tabular}
\end{table}

\paragraph{Real traces (headline).} On Redbench~\cite{redbench2025}---$30$ workloads synthesized from Amazon Redshift's Redset production traces---the exact-query repetition rate \emph{is} our $S(k)$. The empirical savings $1-u_k/k$ spans $0$--$98\%$ (mean $53\%$) with fitted Zipf $\alpha$ from $0$ to ${\approx}3$ (mean $0.84$), a range our synthetic sweep ($\alpha\in\{0,\dots,3\}$) covers. Forecasting the full distinct-query count from only the \emph{first half} of each timeline, a smoothed Good--Toulmin estimator cuts the plug-in's error by $45\%$ ($0.22\!\to\!0.12$ in $u_k/k$). This---not the in-sample fit---is the genuine ahead-of-time forecast of the unexecuted remainder, confirmed on production-derived data. Redbench is \emph{not} a DP benchmark; we use only its repetition structure to externally validate the occupancy forecast.

\paragraph{Calibration and allocation.} Two attacks confirm the per-release noise is correctly calibrated: a single-query membership-inference AUC stays at or below the $e^\eps/(1+e^\eps)$ optimal-\emph{accuracy} reference (within one percentage point for $\eps\le1$), and a differencing reconstruction error scales as $1.5/\eps$. For allocation, the safe $\eps_q{=}B/m$ sizing answers $100\%$ of queries at fresh-release MAE $2.0$ versus an $\eps$-greedy bandit's $3.6$---it \emph{reaches} the bandit's best operating point at zero exploration cost (in the $m\le k/4$ regime tested). Against reactive caches we cite their own figures (Turbo $1.7$--$15.9\times$; ours $14$--$33\times$): the same order of magnitude---though the metrics differ (Turbo \emph{generalizes across overlapping} queries, we reuse only \emph{exact} repeats), so this is a ballpark placement, not a controlled comparison, and we claim \emph{parity}, not superiority.

\section{Conclusion}
\label{sec:concl}

A closed-form model parameterized by $(\{p_i\},k,B)$ predicts expected DP-budget consumption in repeated aggregate SQL workloads, the realized spend concentrating within $O(\sqrt m)$ of it (Prop.~\ref{prop:conc})---the expected spend, a hard static-regime completion guarantee, and a per-release accuracy verdict---from public structure alone, before any query runs and without touching the data. We are explicit about scope: the synthetic grid is a self-consistency check (the empirical $u_k$ is drawn from the very distribution the formula integrates), the temporal extension is first-order, the GROUP BY guarantee assumes a public key domain, and savings genuinely vanish on workloads that never repeat---which the model predicts rather than hides. The external evidence is the Redbench prefix forecast: a $45\%$ error reduction on production-derived traces the model never saw. Code, the $134$-test middleware, and an interactive demo are public.\footnote{\url{https://github.com/canoztas/cmp653-project}}

\begin{thebibliography}{99}\small
\bibitem{dwork2006calibrating} C.\ Dwork, F.\ McSherry, K.\ Nissim, A.\ Smith. Calibrating Noise to Sensitivity in Private Data Analysis. TCC, 2006.
\bibitem{dinur2003revealing} I.\ Dinur, K.\ Nissim. Revealing Information While Preserving Privacy. PODS, 2003.
\bibitem{mcsherry2009privacy} F.\ McSherry. Privacy Integrated Queries (PINQ). SIGMOD, 2009.
\bibitem{johnson2020chorus} N.\ Johnson, J.\,P.\ Near, J.\,M.\ Hellerstein, D.\ Song. Chorus: a Programming Framework for Building Scalable Differential Privacy Mechanisms. IEEE EuroS\&P, 2020.
\bibitem{kotsogiannis2019privatesql} I.\ Kotsogiannis et al. PrivateSQL: A Differentially Private SQL Query Engine. PVLDB 12(11), 2019.
\bibitem{kostopoulou2023turbo} K.\ Kostopoulou, P.\ Tholoniat, A.\ Cidon, R.\ Geambasu, M.\ L\'ecuyer. Turbo: Effective Caching in Differentially-Private Databases. SOSP, 2023.
\bibitem{mazmudar2023cachedp} M.\ Mazmudar, T.\ Humphries, et al. Cache Me If You Can: Accuracy-Aware Inference Engine for Differentially Private Data Exploration. PVLDB 16(4), 2022.
\bibitem{peng2013pioneer} S.\ Peng, Y.\ Yang, Z.\ Zhang, M.\ Winslett, Y.\ Yu. Query Optimization for Differentially Private Data Management Systems (Pioneer). ICDE, 2013.
\bibitem{mundra2026lapras} P.\ Mundra, A.\ Sealfon, Z.\ Sun, Q.\,C.\ Liu. LAPRAS: Learning-Augmented Private Answering for Linear Query Streams. arXiv:2605.01960, 2026.
\bibitem{ge2019apex} C.\ Ge, X.\ He, I.\,F.\ Ilyas, A.\ Machanavajjhala. APEx: Accuracy-Aware Differentially Private Data Exploration. SIGMOD, 2019.
\bibitem{pujol2023multianalyst} D.\ Pujol, A.\ Sun, B.\ Fain, A.\ Machanavajjhala. Multi-Analyst Differential Privacy for Online Query Answering. PVLDB 16(4), 2022.
\bibitem{mironov2017renyi} I.\ Mironov. R\'enyi Differential Privacy. CSF, 2017.
\bibitem{bun2016concentrated} M.\ Bun, T.\ Steinke. Concentrated Differential Privacy: Simplifications, Extensions, and Lower Bounds. TCC, 2016.
\bibitem{flajolet1992} P.\ Flajolet, D.\ Gardy, L.\ Thimonier. Birthday Paradox, Coupon Collectors, Caching Algorithms and Self-Organizing Search. Discrete Applied Math 39(3), 1992.
\bibitem{joagdev1983} K.\ Joag-Dev, F.\ Proschan. Negative Association of Random Variables, with Applications. \emph{Ann.\ Statist.} 11(1), 1983.
\bibitem{li2015matrix} C.\ Li, G.\ Miklau, M.\ Hay, A.\ McGregor, V.\ Rastogi. The Matrix Mechanism: Optimizing Linear Counting Queries under Differential Privacy. VLDB Journal 24(6), 2015.
\bibitem{mckenna2018hdmm} R.\ McKenna, G.\ Miklau, M.\ Hay, A.\ Machanavajjhala. Optimizing Error of High-Dimensional Statistical Queries under Differential Privacy. PVLDB 11(10), 2018.
\bibitem{luo2021dpf} T.\ Luo, M.\ Pan, P.\ Tholoniat, A.\ Cidon, R.\ Geambasu, M.\ L\'ecuyer. Privacy Budget Scheduling. OSDI, 2021.
\bibitem{kuchler2024cohere} N.\ K\"uchler, E.\ Opel, H.\ Lycklama, A.\ Viand, A.\ Hithnawi. Cohere: Managing Differential Privacy in Large Scale Systems. IEEE S\&P, 2024.
\bibitem{goodtoulmin1956} I.\,J.\ Good, G.\,H.\ Toulmin. The Number of New Species, and the Increase in Population Coverage, when a Sample is Increased. Biometrika 43, 1956.
\bibitem{orlitsky2016unseen} A.\ Orlitsky, A.\,T.\ Suresh, Y.\ Wu. Optimal Prediction of the Number of Unseen Species. PNAS 113(47), 2016.
\bibitem{acharya2018inspectre} J.\ Acharya, G.\ Kamath, Z.\ Sun, H.\ Zhang. INSPECTRE: Privately Estimating the Unseen. ICML, 2018.
\bibitem{redbench2025} S.\ Krid, M.\ Stoian, A.\ Kipf. Redbench: A Benchmark Reflecting Real Workloads. aiDM @ SIGMOD, 2025. arXiv:2506.12488.
\end{thebibliography}

\end{document}
